
\section{Related Work}

\paragraph{Prompt injection and tool-integrated applications.}
Indirect prompt injection showed that untrusted data can steer tool-integrated LLM applications \citep{greshake2023indirect}. Subsequent benchmark work formalizes prompt injection attacks and defenses \citep{liu2023promptinjection}. FlowFence-Lite addresses a complementary runtime question: once data is moving through a multi-agent system, how should shared memory, workspaces, and tools be mediated?

\paragraph{Agent security benchmarks.}
AgentDojo provides a dynamic environment for evaluating attacks and defenses in tool-using agents \citep{debenedetti2024agentdojo}. ASB formalizes and benchmarks attacks and defenses in LLM-based agents \citep{zhang2025asb}. Our benchmark is narrower but focuses on multi-agent propagation through shared runtime state and on privacy-specific metrics such as raw leakage, external leakage, cascade size, and privilege reach.

\paragraph{Memory poisoning and agent privacy.}
AgentPoison demonstrates that poisoned memory or knowledge bases can compromise LLM agents \citep{chen2024agentpoison}. MEXTRA and related work highlight privacy risks in LLM agent memory \citep{wang2025mextra}. Our P0 comparator uses an adapted retrieval-memory axis, while the P1 benchmark extends containment to shared memory, shared workspaces, and inter-agent handoffs.

\paragraph{Multi-agent propagation and topology-aware defense.}
Agent Smith illustrates that local compromise can propagate across interacting agents \citep{gu2024agentsmith}. G-Safeguard studies topology-aware guardrails for LLM-based multi-agent systems \citep{wang2025gsafeguard}. FlowFence-Lite similarly treats topology as security-relevant, but focuses on policy-aware privacy containment and safe-view governance of shared artifacts.

\paragraph{Privacy minimization, memory systems, and least privilege.}
AgentDAM evaluates privacy leakage in autonomous web agents \citep{zharmagambetov2025agentdam}, while A-MEM studies memory mechanisms for agents \citep{xu2025amem}. FlowFence-Lite is closer to an information-flow control layer: it combines least-privilege intuitions \citep{saltzer1975protection} and information-flow policy ideas \citep{denning1976lattice} with practical runtime rewriting and quarantine.
